% !PRINTED 392
\weberpage{392}

%<<CONT
solches System $S'$ über, und wir nennen ein solches System
\gesperrt{reell}, wenn $S$ mit $S'$ identisch ist. Ein reelles Siebener-System
enthält daher zu jeder in ihm vorkommenden Wurzel $\xi$ die conjugirte
$\xi'$, und muss folglich wenigstens eine, immer aber eine
ungerade Anzahl von reellen Doppeltangenten enthalten.

Wir leiten jetzt der Reihe nach die Hauptsätze ab:

\gesperrt{1. Es können nicht alle Doppeltangenten imaginär sein.}

Es sei nämlich $\xi, \xi'$ ein conjugirtes Paar und $\eta$ eine dritte
imaginäre Doppeltangente, so dass die drei $\xi, \xi', \eta$ \gesperrt{azygetisch}
sind. (Dies wäre sicher möglich, wenn alle Doppeltangenten
imaginär wären.) Ist $\eta'$ zu $\eta$ conjugirt, so haben wir die beiden
azygetischen Tripel:

\[
\xi,\ \xi',\ \eta;\quad \xi,\ \xi',\ \eta'.
\]

Da hiernach in dem durch das Paar $\xi\eta$ bestimmten \gesperrt{Steiner}'schen
Complexe $\xi'$ nicht vorkommt, so ist dieser Complex
imaginär.

Wir stellen ihn mit seinem conjugirten Complexe zusammen:

\weq{1}{
\begin{aligned}
&1)\quad \xi\eta,\ \xi_1\eta_1,\ \xi_2\eta_2,\ \xi_3\eta_3,\ \xi_4\eta_4,\ \xi_5\eta_5\\
&2)\quad \xi'\eta',\ \xi_1'\eta_1',\ \xi_2'\eta_2',\ \xi_3'\eta_3',\ \xi_4'\eta_4',\ \xi_5'\eta_5',
\end{aligned}
}

indem wir unter $\xi_i', \eta_i'$ die zu $\xi_i, \eta_i$ conjugirten Elemente verstehen,
so dass, falls $\xi_i$ reell ist, $\xi_i = \xi_i'$ zu setzen ist, und umgekehrt
auch aus $\xi_i = \xi_i'$ folgt, dass $\xi_i$ reell ist.

Es sind nun zwei Möglichkeiten zu unterscheiden. Wenn
\gesperrt{erstens} die Complexe 1), 2) ein syzygetisches Paar bilden (\S.\,97),
so haben sie vier syzygetische Elemente gemein. Darunter
können $\xi, \eta, \xi', \eta'$ nicht vorkommen, und wir beschränken daher
die Allgemeinheit nicht, wenn wir annehmen, es seien $\xi_1, \eta_1, \xi_2, \eta_2$
die gemeinschaftlichen Elemente. Diese können in ihrer Gesammtheit
aber nicht verschieden sein von $\xi_1', \eta_1', \xi_2', \eta_2'$, da der
Uebergang zu den accentuirten Buchstaben, d.\,h. zu den conjugirt
imaginären Grössen, wodurch 1) in 2) übergeht, überall
gestattet ist.

Nun ist die Annahme $\xi_1 = \eta_1'$ ausgeschlossen, weil daraus
$\xi_1' = \eta_1$ folgen würde und 1) von 2) nicht verschieden wäre.

Ist $\xi_1 = \xi_1'$, so ist $\xi_1$ \gesperrt{reell}.

Ist aber $\xi_1 = \xi_2'$, $\xi_2 = \xi_1'$,
so ist $\eta_1 = \eta_1'$, also $\eta_1$ \gesperrt{reell},
und davon ist die Annahme $\xi_1 = \eta_2'$, $\xi_1' = \eta_2$, $\xi_2 = \xi_2'$ nicht
wesentlich verschieden.
